\xchapter{Cathedra Sempiterna}

[Reprinted in \emph{My Campaign in Ireland},
printed for private circulation by A. King \& Co., 1896. Put
together from passages taken from Discourse I, \emph{Discourses on the
scope and nature of university education}, pp. 22, 25-28, 1852—\textbf{NR}.]

DEEPLY
do I feel, ever will I protest, for I can appeal to the ample
testimony of history to bear me out, that, in questions of right and
wrong, there is nothing really strong in the whole world, nothing
decisive and operative, but the voice of him, to whom have been
committed the keys of the kingdom and the oversight of Christ's flock.
The voice of Peter is now, as it ever has been, a real authority,
infallible when it teaches, prosperous when it commands, ever taking
the lead wisely and distinctly in its own province, adding certainty
to what is probable, and persuasion to what is certain. Before it
speaks, the most saintly may mistake; and after it has spoken, the
most gifted must obey.

Peter is no recluse, no abstracted student, no
dreamer about the past, no doter upon the dead and gone, no projector
of the visionary. Peter for eighteen hundred years has lived in the
world; he has seen all fortunes, he has encountered all adversaries,
he has shaped himself for all emergencies. If there ever was a power
on earth who had an eye for the times, who has confined himself to the
practicable, and has been happy in his anticipations, whose
words have been deeds, and whose commands prophecies, such is he in
the history of ages, who sits from generation to generation in the
Chair of the Apostles, as the Vicar of Christ and Doctor of His
Church.

It was said by an old philosopher, who declined
to reply to an emperor's arguments, ``It is not safe controverting with
the master of twenty legions.'' What Augustus had in the temporal
order, that, and much more, has Peter in the spiritual. When was he
ever unequal to the occasion? When has he not risen with the crisis?
What dangers have ever daunted him? What sophistry foiled him? What
uncertainties misled him? When did ever any power go to war with
Peter, material or moral, civilized or savage, and got the better?
When did the whole world ever band together against him solitary, and
not find him too many for it?

All who take part with Peter are on the winning
side. The Apostle of Christ says not in order to unsay; for he has
inherited that word which is with power. From the first he has looked
through the wide world, of which he has the burden; and according to
the need of the day and the inspirations of his Lord, he has set
himself, now to one thing, now to another, but to all in season and to
nothing in vain. He came first upon an age of refinement and luxury
like our own; and in spite of the persecutor, fertile in the resources
of his cruelty, he soon gathered, out of all classes of society, the
slave, the soldier, the high-born lady, and the sophist, to form a
people for his Master's honour. The savage hordes came down in
torrents from the north, hideous even to look upon; and Peter went out
with holy water and with benison, and by his very eye he sobered them
and backed them in full career. They turned aside and flooded the
whole earth, but only to be more surely civilized by him, and to be
made ten times more his children even than the older populations they
had overwhelmed. Lawless kings arose, sagacious as the Roman,
passionate as the Hun, yet in him they found their match, and were
shattered, and he lived on. The gates of the earth were opened to the
east and west, and men poured out to take possession; and he and his
went with them, swept along by zeal and charity, as far as they by
enterprise, covetousness, or ambition. Has he failed in his
enterprises up to this hour? Did he, in our fathers' day, fail in his
struggle with Joseph of Germany and his confederates—with Napoleon,
a greater name, and his dependent kings—that, though in another kind
of fight, he should fail in ours? What grey hairs are on the head of
Judah, whose youth is renewed as the eagle's, whose feet are like the
feet of harts, and underneath the Everlasting Arms?

``Thus saith the Lord that created thee, O Jacob,
and formed thee, O Israel. Fear not, for I have redeemed thee, and
called thee by thy name! Thou art Mine.

``When thou shalt pass through the waters, I will
be with thee, and the rivers shall not cover thee.

``When thou shalt walk in the fire, thou shalt not
be burned, and the flame shall not kindle against thee.

``For I am the Lord thy God, the Holy One of
Israel, thy Saviour.

``Fear not, for I am with thee, I am the first,
and I am the last, and besides Me there is no God.''

JOHN HENRY
NEWMAN.

———————

It is not altogether irrelevant to mention here
that in January, 1856, Dr. Newman, having occasion to go to Rome on
business of very great anxiety, he at once, on alighting from the
diligence, went with Father St. John to make a visit of devotion to
the shrine of St. Peter, going there the whole way barefoot. The time
was the middle of the day, when, as was the case in those years, the
streets were very empty, and thus, and screened by his large Roman
cloak, he was able to do so unrecognized and unnoticed—nor was it
ever known except to Father St. John and another.

His friend Dr. Clifford (the Hon. William J. H.
Clifford, late Bishop of Clifton), who with his father Lord Clifford,
had travelled with him from Siena, and with whom he dined that day in
Rome, knew nothing of this until it was mentioned to him on occasion
of his preaching the Cardinal's funeral sermon in 1890.
